\chapter*{Introduction générale}
\phantomsection
\addcontentsline{toc}{chapter}{Introduction générale}
\markboth{Introduction générale}{Introduction générale}

Dans un environnement où les systèmes d'information deviennent le cœur stratégique des organisations, leur capacité à répondre aux besoins métiers de manière intégrée, évolutive et sécurisée constitue un enjeu majeur de performance et de compétitivité. Le secteur des services financiers, en particulier, est marqué par la digitalisation des services bancaires \cite{vives_digital_banking}. Dans l'UEMOA, cette évolution s'inscrit également dans la transformation numérique des services financiers portée par les FinTech \cite{bceao_fintech}, avec des attentes accrues en matière d'instantanéité des transactions et d'expérience utilisateur unifiée.

Cependant, plusieurs organisations utilisent encore des applications séparées qui communiquent peu entre elles. Les données sont alors dispersées, les traitements deviennent plus difficiles à suivre et la qualité du service peut être affectée. Dans les établissements bancaires et les sociétés de gestion, cette difficulté se manifeste notamment lorsque la relation client et les services d'investissement ne sont pas suffisamment reliés.

C'est dans ce contexte que s'inscrit le présent mémoire, réalisé au sein de l'entreprise INNOV4AFRICA dans le cadre d'un stage de deuxième cycle. Le travail effectué porte sur la conception et le développement d'un système d'information bancaire intelligent dénommé \textbf{I-SIB}, avec un focus sur deux modules complémentaires : un module dédié à la gestion de la relation client, inspiré des principes du CRM \cite{payne_frow_crm}, et un module SGI, relatif aux Sociétés de Gestion et d'Intermédiation \cite{crepmf_sgi}.



La suite de ce document est structurée en quatre chapitres.

\begin{itemize}
    \item Le \textbf{premier chapitre} présente l'entreprise d'accueil, expose le contexte spécifique du projet, la problématique identifiée, les objectifs poursuivis ainsi que la démarche méthodologique adoptée.
    \item Le \textbf{deuxième chapitre} est consacré à l'analyse fonctionnelle et non fonctionnelle des besoins.
    \item Le \textbf{troisième chapitre} détaille la conception de la solution, son architecture et les choix technologiques retenus.
    \item Le \textbf{quatrième chapitre} présente les résultats obtenus et propose une analyse critique de la solution développée.
\end{itemize}

La conclusion générale revient sur les apports du projet, les difficultés rencontrées et les perspectives d'amélioration.
